% The pictures of README.md: results of the showcase as the package draws them, one panel per page (each page becomes
% Documentation/readme_<n>.svg). Built by make_readme_figure.sh, which first copies the figures of Showcase/figures
% (run Showcase/make_figures.wl first) to \figdir with one unit of length, heavier lines and larger labels for all of
% them; the formulas are those the showcase computes (\formula). One panel, one theme: the outputs of the first example
% (the electron gas, then phi^4 theory), then exact equations, screening, pairing, field theory. The outputs of the first example are drawn as DiagramPlot draws
% them, whole, one band per order (as the showcase prints a sum order by order). Every panel has the same width, so that
% the README shows them all at the same scale; the titles are captions in README.md. Compiled with \socialpreviewpanel
% defined, the document is instead the one panel of the social preview (social_preview.tex).
\documentclass[multi=panel,border=11pt]{standalone}
\usepackage[T1]{fontenc}
\usepackage[sc]{mathpazo}
\usepackage{amsmath,tikz,xcolor,array,varwidth,environ}
\usetikzlibrary{decorations.pathmorphing}
\providecommand{\figdir}{../Showcase/figures}
\newcommand{\fig}[1]{\input{\figdir/#1.tex}}
% a formula written by ExpressionToLaTeX in the showcase (figures/formula_NAME.tex), so that every number is computed
\providecommand{\formuladir}{../Showcase/figures}
\newcommand{\formula}[1]{\input{\formuladir/formula_#1.tex}\unskip}
\newcommand{\diagramgap}{\hskip 0.6em\relax}
% lines of diagrams in a paragraph kept apart, so that the labels of two lines never touch
\newcommand{\figlines}{\lineskiplimit=7pt\lineskip=7pt\relax}
% the rest of a series, kept on the line of the last diagram
\newcommand{\rest}{\unskip\nobreak\diagramgap\mbox{${}+\cdots$}}
% a panel: rows "left-hand side = diagrams", the diagrams as a ragged paragraph in the rest of the width, the whole
% centred in the panel (the left-hand sides are measured first, without the diagrams)
\newlength{\panelwidth}\setlength{\panelwidth}{10.4cm}
\newlength{\lhswidth}
\NewEnviron{rows}{%
  \settowidth{\lhswidth}{\begin{tabular}{@{}r@{}>{\setbox0\vbox\bgroup\arraybackslash}l<{\egroup}@{}}\BODY\end{tabular}}%
  \makebox[\panelwidth][c]{\begin{tabular}{@{}r@{}>{\begin{varwidth}[t]{\dimexpr\panelwidth-\lhswidth\relax}%
    \raggedright\arraybackslash\figlines}l<{\end{varwidth}}@{}}\BODY\end{tabular}}}
% the space between two rows (\\[...] would add nothing after a row of two lines)
\newcommand{\rowgap}{\\\noalign{\vskip 11pt}}
% a sum written order by order, as the showcase prints it: \ordersum{lhs}{\orderrow{order}{diagrams} ...}, the
% left-hand side on the left, right-aligned at the width of the widest one of the panel (the argument of orderpanel), and
% on its right one band per order, the order on the left of the band and the diagrams as a ragged paragraph
\definecolor{bandA}{HTML}{EEF3F9}
\definecolor{bandB}{HTML}{F7F9FC}
\definecolor{accent}{HTML}{1F4E79}
\newcounter{band}
\newlength{\orderlabelwidth}\setlength{\orderlabelwidth}{0.6cm}
\newlength{\bandwidth}
\newenvironment{orderpanel}[1]{\settowidth{\lhswidth}{$\displaystyle #1={}$}%
  \setlength{\bandwidth}{\dimexpr\panelwidth-\lhswidth\relax}\begin{minipage}{\panelwidth}\parindent=0pt\relax}
  {\end{minipage}}
\newcommand{\ordersum}[2]{\par\noindent\makebox[\lhswidth][r]{$\displaystyle #1={}$}%
  \parbox[t]{\bandwidth}{\setcounter{band}{0}\offinterlineskip #2}\par}
\newcommand{\orderrow}[2]{\stepcounter{band}%
  \ifodd\value{band}\def\bandcolor{bandA}\else\def\bandcolor{bandB}\fi
  \par\noindent{\setlength{\fboxsep}{4pt}\colorbox{\bandcolor}{\parbox[t]{\dimexpr\bandwidth-2\fboxsep\relax}{%
  \rightskip=0pt plus 1fil\parfillskip=0pt plus 1fil\baselineskip=0pt\figlines\let\diagramgap\banddiagramgap
  \hangindent=\orderlabelwidth\hangafter=1\noindent\makebox[\orderlabelwidth][l]{\small\color{accent}$#1$}#2\par}}}\par}
% between two diagrams of a band, the space at the end of the line before it dropped
\newcommand{\banddiagramgap}{\unskip\hskip 0.6em\relax}
% the space between two sums of a panel
\newcommand{\ordergap}{\par\vspace{11pt}}
\begin{document}
\ifdefined\socialpreviewpanel
% the panel of the social preview: Phi[G] and its derivatives up to second order, the rest as + ... (a picture, not the
% output of a command)
\begin{panel}
\begin{rows}
$\Phi[G]={}$ & \fig{luttinger_ward_1}\diagramgap\fig{luttinger_ward_2}\rest\rowgap
$\Sigma=\dfrac{\delta\Phi}{\delta G}={}$ & \fig{self_energy_1}\diagramgap\fig{self_energy_2}\rest\rowgap
$\Gamma^0=\dfrac{\delta\Sigma}{\delta G}={}$ & \fig{irreducible_vertex_1}\diagramgap\fig{irreducible_vertex_2}\rest
\end{rows}
\end{panel}
\else
% 1: the first example of README.md: the Green function to second order
\begin{panel}
\begin{orderpanel}{G(x,y)}
\ordersum{G(x,y)}{%
\orderrow{V^0}{\fig{green_function_0}}
\orderrow{V}{\fig{green_function_1}}
\orderrow{V^2}{\fig{green_function_2}}}
\end{orderpanel}
\end{panel}
% 2: the first example of README.md: Phi[G] and its derivatives to second order
\begin{panel}
\begin{orderpanel}{\Gamma^0=\dfrac{\delta\Sigma}{\delta G}}
\ordersum{\Phi[G]}{%
\orderrow{V}{\fig{luttinger_ward_1}}
\orderrow{V^2}{\fig{luttinger_ward_2}}}\ordergap
\ordersum{\Sigma=\dfrac{\delta\Phi}{\delta G}}{%
\orderrow{V}{\fig{self_energy_1}}
\orderrow{V^2}{\fig{self_energy_2}}}\ordergap
\ordersum{\Gamma^0=\dfrac{\delta\Sigma}{\delta G}}{%
\orderrow{V}{\fig{irreducible_vertex_1}}
\orderrow{V^2}{\fig{irreducible_vertex_2}}}
\end{orderpanel}
\end{panel}
% 3: the first example of README.md: the propagator of phi^4 theory to second order
\begin{panel}
\begin{orderpanel}{\langle\phi_a\phi_b\rangle}
\ordersum{\langle\phi_a\phi_b\rangle}{%
\orderrow{\lambda^0}{\fig{phi4_propagator_0}}
\orderrow{\lambda}{\fig{phi4_propagator_1}}
\orderrow{\lambda^2}{\fig{phi4_propagator_2}}}
\end{orderpanel}
\end{panel}
% 4: exact equations of many-body theory, drawn with blobs: Dyson, Bethe-Salpeter, the vertex of the density, parquet
\begin{panel}
\begin{rows}
\fig{exact_dyson_lhs}\,${}={}$ & \fig{exact_dyson}\rowgap
\fig{exact_vertex_lhs}\,${}={}$ & \fig{exact_bethe_salpeter}\rowgap
\fig{exact_lambda_lhs}\,${}={}$ & \fig{exact_vertex_equation}\rowgap
\fig{exact_vertex_lhs}\,${}={}$ & \fig{exact_parquet}
\end{rows}
\end{panel}
% 5: the screened interaction: Hedin's equations, and the functional Phi[G, W] whose derivatives give GW
\begin{panel}
\begin{rows}
$\Sigma(x,y)={}$ & \fig{hedin_sigma}\rowgap
\fig{hedin_w_lhs}\,${}={}$ & \fig{hedin_w}\rowgap
\fig{hedin_p_lhs}\,${}={}$ & \fig{hedin_p}\rowgap
\fig{hedin_lambda_lhs}\,${}={}$ & \fig{hedin_lambda}\rowgap
$\Phi[G,W]={}$ & \fig{phi_gw_1}\diagramgap\fig{phi_gw_2}\rest\rowgap
$\Sigma=\dfrac{\delta\Phi}{\delta G}={}$ & \fig{sigma_gw_1}\rest\rowgap
$P=-2\dfrac{\delta\Phi}{\delta W}={}$ & \fig{polarization_gw_0}\diagramgap\fig{polarization_gw_1}\rest
\end{rows}
\end{panel}
% 6: pairing: the superfield of Nambu and Gor'kov
\begin{panel}
\begin{rows}
$\Phi[G,\mathcal G_{\psi\psi},\mathcal G_{\bar\psi\bar\psi}]={}$ & \fig{nambu_phi_components_1}\rest\rowgap
\fig{gorkov_g_lhs}\,${}={}$ & \fig{gorkov_g}\rowgap
\fig{gorkov_gpp_lhs}\,${}={}$ & \fig{gorkov_gpp}\rowgap
\fig{nambu_sigma_lhs}\,${}={}$ & \fig{nambu_hartree_fock}\rowgap
\fig{nambu_gap_lhs}\,${}={}$ & \fig{nambu_gap}
\end{rows}
\end{panel}
% 7: field theory: phi^4 in dimensional regularization to two loops, QED at one loop (each 1PI function with its main
% result), the gluon self-energy of QCD; the formulas are those the showcase computes
\begin{panel}
\begin{rows}
$\Sigma^{(2)}(p)={}$ & \fig{phi4_self_energy_2_alone}$\displaystyle{}=\formula{dimreg_sunset_expanded}$\rowgap
$\beta(\lambda)={}$ & $\displaystyle\formula{dimreg_beta},\qquad \kappa=\frac1{16\pi^2}$\rowgap
$-i\Sigma_2={}$ & \fig{qed_self_energy}\qquad$\delta_2=\displaystyle\formula{qed_delta2}$\rowgap
$i\Pi_2^{\mu\nu}={}$ & \fig{qed_vacuum_polarization}\qquad$\beta(e)=\displaystyle\formula{qed_beta}$\rowgap
$ie\,\delta\Gamma^\mu={}$ & \fig{qed_vertex}\qquad$F_2(0)=\displaystyle\formula{qed_F2_zero}$\rowgap
$\Pi_{\rm QCD}(x,y)={}$ & \fig{qcd_gluon_self_energy}
\end{rows}
\end{panel}
\fi
\end{document}
